\documentclass[11pt,a4paper]{article}

\usepackage[T1]{fontenc}
\usepackage[utf8]{inputenc}
\usepackage{lmodern}
\usepackage[a4paper,margin=27mm]{geometry}
\usepackage{amsmath,amssymb,amsthm,mathtools}
\usepackage{booktabs,array}
\usepackage{microtype}
\usepackage{xcolor}
\usepackage{enumitem}
\usepackage{url}
\usepackage[colorlinks=true,linkcolor=blue!55!black,citecolor=blue!55!black,
            urlcolor=blue!55!black]{hyperref}

\newtheorem{theorem}{Theorem}
\newtheorem{proposition}[theorem]{Proposition}
\newtheorem{lemma}[theorem]{Lemma}
\newtheorem{remark}[theorem]{Remark}
\newcommand{\Z}{\mathbb{Z}}
\newcommand{\R}{\mathcal{R}}
\newcommand{\T}{\mathsf{T}}
\newcommand{\E}{\mathbb{E}}
\newcommand{\Prb}{\mathop{\rm Pr}}
\newcommand{\HighBits}{\mathsf{HighBits}}
\newcommand{\LowBits}{\mathsf{LowBits}}
\newcommand{\pk}{\mathsf{pk}}
\newcommand{\sk}{\mathsf{sk}}
\newcommand{\Sign}{\mathsf{Sign}}
\newcommand{\Verify}{\mathsf{Verify}}
\newcommand{\KeyGen}{\mathsf{KeyGen}}
\newcommand{\norm}[1]{\left\lVert #1\right\rVert}

\title{\bfseries Global-Sign Correlation Breaks BiT-128\\
\large Practical Equivalent-Key Recovery and Chosen-Message Forgery}
\author{Martin Feussner\\
\small Selmer Center, University of Bergen\\
\small \texttt{martin.feussner@uib.no}}
\date{25 September 2026}

\begin{document}
\maketitle

\begin{center}
\fbox{\parbox{0.91\linewidth}{\small\textbf{Provenance and review status.}
The attack, derivation, implementation, experiments, and initial manuscript
were produced independently by OpenAI Codex using Daybreak Blue at max
reasoning effort during an AI-run audit of NGCC signature candidates.  At the
date above, no human cryptanalyst has yet independently verified the argument,
implementation, or measurements.  This report should therefore be treated as
a preliminary disclosure for human review and reproduction.}}
\end{center}
\vspace{0.7em}

\begin{abstract}
BiT is a Fiat--Shamir-with-aborts lattice signature submitted to the
Next-generation Commercial Cryptographic Algorithms Program.  Its response is
formed as $z=y+rsc$, where a single hidden sign $r\in\{-1,+1\}$ is shared by
all coefficients.  BiT then applies a one-dimensional rejection rule to every
coefficient separately.  We show that this rule produces the advertised
one-dimensional marginals but does not remove the joint dependence created by
the shared sign.  The defect is especially damaging because the response to
the public constant component of the secret is included in every signature.
For BiT-128, its 30 challenge positions identify the hidden sign with about
$96.2\%$ hard-decision accuracy; the posterior mean retains about $88.9\%$ of
the information of a revealed sign.

We turn the leakage into a streaming weighted regression over
$\Z[x]/(x^{256}+1)$.  Negacyclic Fourier diagonalisation recovers the three
ternary secret polynomials from ordinary signatures.  A small search over the
least certain rounded coefficients is checked using the compressed public key.
Once those polynomials are known, public-key compression permits construction
of a short error and low part that form an equivalent signing key.

An end-to-end experiment against the submitted BiT-128 reference
implementation used 200,000 signatures.  The direct estimator left five of 768
secret coefficients incorrect.  Public-key correction found the valid
assignment after 1,596 subset tests.  The recovered secret was completed to a
different valid signing key, and the unmodified submitted verifier accepted a
1,504-byte signature on a fresh message.  The entire single-threaded run,
including reference signing, used 139 CPU-seconds on a 3.2 GHz Intel Core
i5-6500 and negligible attack memory.  The attack is information-theoretic,
does not exploit the hash function, and remains valid in the random-oracle
model.  The fixed 200,000-query procedure forged on three of four independently
generated keys; the fourth forged after extending the transcript to 250,000
signatures.
\end{abstract}

\section{Introduction}

Fiat--Shamir-with-aborts signatures use rejection sampling to make a response
statistically independent of the signing key.  That independence is a joint
distribution requirement: an adversary receives a complete vector, not a
collection of isolated scalar samples.  Replacing a vector rejection ratio by
the product of ratios for its scalar marginals is generally invalid when a
latent random variable is shared by all coordinates.

BiT (``Bimodal Triangular distribution based lattice signatures'') is an NGCC
digital-signature candidate built around exactly such a latent variable
\cite{BiT}.  It samples one bit $b$, writes $r=(-1)^b$, and forms
\[
             z = y+rsc.
\]
The mask coefficients are independent triangular samples, but the sign $r$ is
global.  BiT's rejection procedure derives the correct scalar mixture for each
coordinate and accepts each coordinate with the corresponding scalar ratio.
This proves a marginal statement.  It does not transform the vector mixture
into a secret-independent product distribution.

The distinction is directly exploitable.  BiT's extended secret is
$s=(1,s_0,e)$.  Its first response polynomial therefore contains
$z_{1,0}=y_{1,0}+rc$.  At each nonzero challenge position the offset is the
known value $r$, so these public response coefficients act as repeated noisy
observations of the global sign.  Once a posterior for $r$ is available, the
remaining public response polynomials
$z_{1,1}=y_{1,1}+r(s_0c)$ become noisy linear observations of $s_0$.

Our contributions are as follows.

\begin{itemize}[leftmargin=1.5em]
  \item We identify the joint-distribution error in BiT's coordinatewise
        bimodal rejection sampler and give an explicit expression for the
        residual dependence.
  \item We derive a sign distinguisher from $z_{1,0}$.  For the BiT-128
        parameters it has about $96.2\%$ hard-decision accuracy from one
        signature.
  \item We give a streaming key-recovery algorithm based on posterior-weighted
        regression.  A negacyclic Fourier transform reduces recovery to 768
        scalar estimates and requires only $O(N)$ stored state per secret
        polynomial.
  \item We show how to validate and correct uncertain coefficients using the
        compressed public key, and how to complete recovered $s_0$ to an
        equivalent signing key without recovering the original error $e$.
  \item We demonstrate an accepted fresh-message forgery against the submitted
        BiT-128 implementation using 200,000 signatures and about 2.3 minutes
        of single-threaded CPU time on an ordinary 2015 desktop processor.
\end{itemize}

This is a full key-recovery and forgery attack on the claimed 128-bit instance.
It is not a distinguisher, a side channel, or a reduction in an asymptotic
security estimate.  The attack needs no control over messages beyond obtaining
distinct valid signatures; it can therefore be mounted passively when enough
known-message signatures are available.

\section{The relevant part of BiT}

We describe only the components needed for the attack.  The notation and
parameter names follow the submission specification \cite{BiT}.  Let
\[
  \R_q=\Z_q[x]/(x^N+1), \qquad N=256, \qquad q=26881
\]
for BiT-128.  Multiplication by a challenge polynomial $c$ is a negacyclic
linear map, which we denote by $C$ when using coefficient vectors.

\subsection{Keys and public-key compression}

The secret consists of $s_0\in S_1^m$ and $e\in S_1^n$, where every
coefficient lies in $\{-1,0,1\}$ and $(m,n)=(3,3)$.  With public matrix
$A_0\in\R_q^{n\times m}$, key generation computes
\[
  u=A_0s_0+e,
  \qquad
  (b_1,b_0)=\bigl(\HighBits(u,\gamma_b),
                         \LowBits(u,\gamma_b)\bigr),
  \qquad \gamma_b=16.
\]
The public key contains the seed defining $A_0$ and $b_1$; the signing key also
contains $s_0,e,b_0$.  Defining
\[
 A=\bigl[\widehat q j_n-\gamma_b b_1,\ A_0,\ I_n\bigr],
 \qquad s=(1,s_0^T,e^T)^T,
\]
where $\widehat q=2^{-1}\bmod q$, gives
$As=\widehat qj_n+b_0$.

\subsection{Signing and the shared sign}

For a message, the signer samples independent triangular masks
\[
 y=(y_{1,0},y_{1,1},y_2)
   \leftarrow
 \T_{\gamma_{1,0}}\times
 \T_{\gamma_{1,1}}^m\times
 \T_{\gamma_{1,2}}^n.
\]
It hashes the high part of $Ay$ and the message representative to obtain a
challenge $c$ of Hamming weight $\tau$.  The submitted BiT challenge has
coefficients in $\{0,1\}$, with exactly $\tau=30$ ones for BiT-128.  One bit
$b$ is sampled and the global sign $r=(-1)^b$ is applied to all components:
\begin{align*}
 z_{1,0}&=y_{1,0}+rc,\\
 z_{1,1}&=y_{1,1}+r(s_0c),\\
 z_2&=y_2+r(ec).
\end{align*}
The signer applies coefficientwise rejection, performs further high-bits and
hint checks, and returns $(z_1,h,\widetilde c)$, where
$z_1=(z_{1,0},z_{1,1})$.  The $z_2$ component is compressed into the hint $h$.

Table~\ref{tab:params} lists the parameters that control the attack.

\begin{table}[ht]
\centering
\caption{Relevant submitted BiT-128 parameters.}
\label{tab:params}
\begin{tabular}{@{}lll@{}}
\toprule
Parameter & Value & Role \\
\midrule
$N,q$ & $256,26881$ & ring degree and modulus \\
$(n,m)$ & $(3,3)$ & module dimensions \\
$\eta$ & $1$ & ternary secret bound \\
$\tau$ & $30$ & challenge weight \\
$\gamma_{1,0}$ & $8$ & mask for the known constant component \\
$\gamma_{1,1}$ & $1024$ & mask for $s_0c$ \\
$\gamma_{1,2}$ & $2048$ & mask for $ec$ \\
$\gamma_b$ & $16$ & public-key compression interval \\
Public key & 1048 bytes & submitted encoding \\
Signature & 1504 bytes & submitted encoding \\
\bottomrule
\end{tabular}
\end{table}

\subsection{Scalar triangular rejection}

Ignoring a common normalising factor, the triangular mass function is
\[
 f_\gamma(x)=\max\{\gamma-|x|,0\}.
\]
For a scalar secret offset $v$, the source marginal after hiding the sign is
\[
 F_{\gamma,v}(x)=\frac{f_\gamma(x-v)+f_\gamma(x+v)}{2}.
\]
BiT defines a target $G_{\gamma,\beta}(x)$ using a public bound $\beta$ and
accepts a scalar response with probability proportional to
\[
                a_v(x)=\frac{G_{\gamma,\beta}(x)}
                              {F_{\gamma,v}(x)}.
\]
The procedure is then applied separately to every coefficient.  At the scalar
level this is correct: averaging over $r$ transforms $F_{\gamma,v}$ into the
chosen target.  The next section shows why applying these ratios independently
does not hide a sign shared by the entire vector.

\section{The joint-distribution flaw}

Consider $d$ response coordinates with independent masks and secret offsets
$v=(v_1,\ldots,v_d)$, but a single sign $r$.  Before rejection, the joint
source is
\begin{equation}
 P_v(z)=\frac12\prod_{i=1}^d f_i(z_i-v_i)
       +\frac12\prod_{i=1}^d f_i(z_i+v_i).
 \label{eq:true-source}
\end{equation}
It is not the product of scalar mixtures
$\prod_i F_i(z_i)$.  Nevertheless, BiT applies the acceptance product
\begin{equation}
                     a_v(z)=\prod_{i=1}^d\frac{G_i(z_i)}{F_i(z_i)}.
 \label{eq:accept-product}
\end{equation}

For a fixed sign, define the accepted one-coordinate law
\begin{equation}
 q_{i,r}(x)=\frac{f_i(x-rv_i)G_i(x)/F_i(x)}{\mathcal Z_i},
 \qquad
 \mathcal Z_i=\sum_x f_i(x-v_i)G_i(x)/F_i(x).
 \label{eq:qir}
\end{equation}
Symmetry gives the same normaliser for $r=+1$ and $r=-1$, and
$q_{i,-}(x)=q_{i,+}(-x)$.

\begin{proposition}[Marginals are hidden but the joint law is not]
After BiT's coordinatewise rejection, the response law is
\begin{equation}
 Q_v(z)=\frac12\prod_{i=1}^d q_{i,+}(z_i)
       +\frac12\prod_{i=1}^d q_{i,-}(z_i).
 \label{eq:accepted-mixture}
\end{equation}
Every scalar marginal equals the advertised target $G_i/\mathcal Z_i$, but, except in
degenerate cases, the joint distribution depends on $v$.
\end{proposition}

\begin{proof}
Conditioned on $r$, independent masks and independent rejection coins factor,
which gives the two products in \eqref{eq:accepted-mixture}.  Averaging one
coordinate over the two signs gives
\[
 \frac{q_{i,+}(x)+q_{i,-}(x)}2
 =\frac{G_i(x)}{\mathcal Z_i}.
\]
The mixture of two products in \eqref{eq:accepted-mixture} is generally not the
product of these averages.  For example, let $X_i$ denote the $i$th accepted
response coordinate and put $\mu_i(v_i)=\E[X_i\mid r=+1]$.
Conditional independence and symmetry give,
for distinct coordinates,
\[
    \E[X_iX_j]=\mu_i(v_i)\mu_j(v_j),
    \qquad \E[X_i]=\E[X_j]=0.
\]
Thus the shared sign leaves a secret-dependent cross-coordinate covariance
even though both scalar marginals have the intended distribution.
\end{proof}

A correct vector rejection sampler would have to use the true source
\eqref{eq:true-source} in a joint ratio.  Multiplying ratios derived from the
scalar mixtures silently replaces a shared latent sign by independent latent
signs, which is a different distribution.

\section{Recovering the hidden sign from the known component}

BiT makes the joint leak unusually accessible because the first component of
the extended secret is the public constant $1$.  At each challenge position,
\[
                  X=z_{1,0,i}=Y+r,
\]
with mask width $\gamma_{1,0}=8$.  From \eqref{eq:qir}, the scalar likelihood
ratio is
\begin{equation}
 \frac{q_+(x)}{q_-(x)}
   =\frac{f_8(x-1)}{f_8(x+1)}.
 \label{eq:scalar-llr}
\end{equation}
The target mass and acceptance probability cancel.  For the 15 possible
values $x=-7,\ldots,7$, the conditional numerator masses are
\begin{equation}
\begin{aligned}
 63q_+(x)&=(0,1,2,3,4,5,6,7,8,7,6,5,4,3,2),\\
 63q_-(x)&=(2,3,4,5,6,7,8,7,6,5,4,3,2,1,0).
\end{aligned}
\label{eq:small-law}
\end{equation}
In particular, $x=7$ determines $r=+1$ and $x=-7$ determines $r=-1$.

A signature contains $\tau=30$ such samples sharing the same $r$.  Let $S_t$
be the support of its challenge.  The log-likelihood ratio and posterior mean
are
\begin{align}
 \Lambda_t
   &=\sum_{i\in S_t}
       \log\frac{f_8(z_{1,0,i}-1)}{f_8(z_{1,0,i}+1)},
       \label{eq:llr}\\
 w_t&=\E[r_t\mid z_{1,0},c_t]
       =\tanh(\Lambda_t/2).
       \label{eq:posterior}
\end{align}
Numerical convolution of the finite law \eqref{eq:small-law} gives a Bayes
hard-decision success probability of approximately $0.962$.  It also gives
\begin{equation}
                       \E[w_t^2]\approx0.8894.
 \label{eq:w2}
\end{equation}
The latter is the useful quantity for soft-decision recovery.  A value of one
would correspond to a fully revealed sign, while zero would convey no sign
information.  Thus one BiT-128 signature retains almost 89 percent of the
linear-regression information available if $b$ were explicitly published.

The additional BiT rejection and hint conditions do not repair this defect.
In our final-signature measurements the average $w_t^2$ was $0.8896$, matching
\eqref{eq:w2} to four significant figures.  In a separate 20,000-signature
run instrumented only to record the sign after each signature had already been
accepted, the hard decision $\operatorname{sign}(w_t)$ was correct in
$96.265\%$ of signatures and $Q^{-1}\sum_t w_tr_t=0.889743$.  The attack itself
does not use this instrumentation.

\section{Streaming recovery of the ternary secret}

Let $s_j\in\{-1,0,1\}^{N}$ be one of the three polynomials of $s_0$, and let
$C_t$ be negacyclic multiplication by the public challenge $c_t$.  Conditional
on $r_t$, each accepted response coordinate has mean
\begin{equation}
 m(v)=
 \frac{\sum_x x f_{1024}(x-v)G_{1024,30}(x)/F_{1024,v}(x)}
      {\sum_x   f_{1024}(x-v)G_{1024,30}(x)/F_{1024,v}(x)}.
 \label{eq:m-v}
\end{equation}
This is odd in $v$.  On the complete range $|v|\leq30$, it is close to the
identity: $m(1)=0.99987$ and $m(30)=29.1461$.  The convolution values arising
from ternary secrets and weight-30 challenges are concentrated much nearer
zero, where the approximation $m(v)\simeq v$ is tighter.

Using the posterior \eqref{eq:posterior}, the public response model is therefore
\begin{equation}
                 z_{t,j}\simeq w_t C_t s_j+\varepsilon_{t,j}.
 \label{eq:regression}
\end{equation}
The noise has standard deviation about $418$ per coordinate, but each
signature supplies $N=256$ equations for every secret polynomial.  Ordinary
weighted least squares gives
\begin{equation}
 \widetilde s_j=
 \left(\sum_t w_t^2C_t^*C_t\right)^{-1}
 \left(\sum_t w_tC_t^*z_{t,j}\right).
 \label{eq:normal}
\end{equation}

The operators $C_t$ are negacyclic circulants and are diagonal at the roots of
$X^N+1$.  If hats denote this negacyclic Fourier transform, then every
frequency can be estimated independently:
\begin{equation}
 \widehat{\widetilde s}_{j,k}=
 \frac{\sum_t w_t\overline{\widehat c_{t,k}}\widehat z_{t,j,k}}
      {\sum_t w_t^2|\widehat c_{t,k}|^2}.
 \label{eq:fourier-estimator}
\end{equation}
The attack streams signatures through these numerator and denominator sums;
it need not store the signatures.  After an inverse transform, each real
coefficient is rounded to $\{-1,0,1\}$.

For $Q$ signatures, the arithmetic cost is
$O(Q(m+1)N\log N)$ and the stored attack state is $O(mN)$.  These costs are
small beside acquiring signatures.  The attack uses the challenge and response
already present in the public signature and never uses a signing-time
measurement.

\section{Public correction and equivalent-key completion}

At practical query counts, a few estimates may lie on the wrong side of a
ternary rounding boundary.  The public key both identifies those errors and
allows completion to a signing key.

\subsection{Low-confidence subset search}

For an estimate $a$, define its distance to the nearest relevant decision
boundary by
\[
                         d(a)=\bigl||a|-1/2\bigr|.
\]
Sort all 768 estimates by $d(a)$.  For a coefficient rounded to $\pm1$, the
alternate value is zero.  For a coefficient rounded to zero, the alternate is
the adjacent sign indicated by $a$.  We enumerate subsets of the 24 least
certain alternates in Gray-code order.

Each assignment determines $A_0s_0$.  Before performing a full check, three
coefficients are tested against the public condition
\begin{equation}
 \exists e'\in\{-1,0,1\}:\quad
 \HighBits\bigl((A_0s_0)_i+e',16\bigr)=(b_1)_i.
 \label{eq:public-filter}
\end{equation}
For a wrong assignment, one such equality holds with probability roughly the
width of one compression interval divided by $q$; three independent checks
therefore discard essentially every wrong subset.  Gray-code updates require
only one precomputed column contribution per subset.

\subsection{Completing an equivalent signing key}

Suppose $s_0$ is correct.  For each coefficient, choose any
$e'_i\in\{-1,0,1\}$ satisfying \eqref{eq:public-filter} and set
\begin{equation}
 b'_0=\LowBits(A_0s_0+e',16).
 \label{eq:completion}
\end{equation}
Such a choice always exists because the original error is one witness.  It is
usually possible to choose $e'_i=0$; only coefficients near compression
boundaries require $\pm1$.  Then
\[
       A_0s_0+e'=16b_1+b'_0,
\]
so $(s_0,e',b'_0)$ satisfies exactly the key relation used by signing and
verification.  The signing-randomness seed can be chosen afresh, and the public
key hash is recomputed from the public key.  This yields an equivalent secret
key, although $e'$ need not equal the key generator's original $e$.

This completion is also a public validation test for recovery.  An incorrect
ternary vector makes $A_0s_0$ pseudorandomly inconsistent with hundreds of
compressed public coefficients; false acceptance is negligible.

\section{End-to-end experiment}

\subsection{Method}

We evaluated the attack against the submitted BiT-128 reference algorithm with
its default SM3-based hashing and exact submitted encodings.  One key pair was
generated, and the signing oracle was queried on distinct 24-byte messages
containing a counter.  Recovery consumed only the public key and the resulting
signatures.  The original secret was retained outside the attack and used only
after each public validation step to count diagnostic coefficient errors.

The implementation used double-precision radix-2 transforms at the 256 roots
of $X^{256}+1$.  It accumulated \eqref{eq:fourier-estimator}, rounded the
inverse transform, and applied the 24-position public correction described
above.  A fresh equivalent key was packed using \eqref{eq:completion}.  The
forgery message was distinct from every oracle query, and the final verdict was
obtained from the submitted verifier without modification.

The experiment ran single-threaded on an Intel Core i5-6500 at 3.2 GHz with
16 GiB RAM, using GCC 13.3.0 with optimisation enabled.  CPU time was measured
over key recovery and the reference signing calls.  Attack state consists of a
few arrays of 256 complex doubles and is far below one megabyte; the reference
signer and temporary objects also fit easily in ordinary laptop memory.

\subsection{Results}

Table~\ref{tab:main-result} gives the complete 200,000-signature run.  The
direct estimator's coefficient RMSE was $0.186328$.  Only five ternary
coefficients rounded incorrectly.  All five occurred among the 24 positions
selected solely by public confidence scores.  The Gray-code enumeration found
the publicly valid assignment after 1,596 subset tests, far short of its
$2^{24}$ cap.

\begin{table}[ht]
\centering
\caption{End-to-end BiT-128 attack result.}
\label{tab:main-result}
\begin{tabular}{@{}lr@{}}
\toprule
Quantity & Result \\
\midrule
Signature queries & 200,000 \\
Posterior mean $Q^{-1}\sum_tw_t^2$ & 0.889585 \\
Direct-estimate RMSE & 0.186328 \\
Wrong coefficients before correction & 5 of 768 \\
Low-confidence positions searched & 24 \\
Subset tests before public match & 1,596 \\
Wrong coefficients after correction & 0 of 768 \\
Nonzero coefficients in completed $e'$ & 31 of 768 \\
CPU time including signing & 139.3 seconds \\
Forged signature size & 1,504 bytes \\
Fresh-message verifier result & accepted \\
\bottomrule
\end{tabular}
\end{table}

The completed key deliberately preferred $e'_i=0$ whenever possible and had
only 31 nonzero error coefficients.  It was therefore plainly different from
a normally sampled error polynomial, yet it met the public relation and signed
correctly.  The verifier accepted the resulting signature on the unqueried
message.  This demonstrates the complete chain
\[
 \text{signatures}\longrightarrow s_0
 \longrightarrow(s_0,e',b'_0)
 \longrightarrow\text{fresh forgery}.
\]

The measured RMSE followed the expected $Q^{-1/2}$ trend: it was $0.3763$ at
50,000 queries, $0.2743$ at 100,000, $0.2166$ at 150,000, and $0.1863$ at
200,000.  Public correction is therefore doing bounded finishing work rather
than compensating for a nonconvergent estimator.

\subsection{Additional-key checks}

We repeated the fixed 200,000-query procedure on three further keys, without
changing the 24-position correction list.  Table~\ref{tab:replications}
reports all four completed key trials, including the main experiment above.
Three recovered an equivalent key and produced an accepted fresh-message
forgery.  On trial D, at least one of the five remaining errors lay outside the
24 selected positions, so the fixed correction search exhausted its cap.  The
failure is reported rather than discarded; an attacker can respond by
collecting more signatures or enlarging the public correction list.

\begin{table}[ht]
\centering
\caption{Fixed-parameter checks on four independently generated BiT-128 keys.}
\label{tab:replications}
\begin{tabular}{@{}crrrl@{}}
\toprule
Trial & RMSE & Direct errors & Subsets tested & Result \\
\midrule
A & 0.186328 & 5 & 1,596 & accepted forgery \\
B & 0.185175 & 8 & 41,055 & accepted forgery \\
C & 0.185637 & 4 & 16,131 & accepted forgery \\
D & 0.194541 & 5 & 16,777,216 & correction cap reached \\
\bottomrule
\end{tabular}
\end{table}

This small sample is a robustness check, not a precise estimate of attack
success probability.  It shows both that the main result is reproducible on
other keys and that 200,000 queries with a 24-position list is not a
deterministic success threshold.  Continuing trial D to 250,000 signatures
reduced its direct estimate to one error with RMSE $0.173932$; the same
24-position correction then found the exact key after 256 subset tests, and the
fresh-message forgery was accepted.  Thus all four tested keys were broken with
at most 250,000 signatures.

\subsection{Why this is a mathematical attack}

No timing, cache trace, fault, malformed input, or implementation-specific
undefined behaviour is used.  The only inputs to recovery are ordinary public
keys and valid signatures.  Hash outputs are treated as independent random
challenges; the attack does not predict, collide, or invert SM3.  The leak is
present in the specified probability distribution, so replacing the hash or
making the implementation constant-time does not address it.

\section{Other parameter sets}

The same joint-distribution defect is present whenever one global sign is
combined with coordinatewise scalar rejection.  BiT-256 and BiT-512 also
publish the response to the known constant component, so the sign-posterior
step applies directly.  Their parameters use larger masks for $s_0c$, which
raises the number of signatures needed by the simple regression in this paper.
We therefore claim a measured full break only for BiT-128.  The higher sets
should not be considered unaffected: they instantiate the same invalid joint
sampling argument, and improved nonlinear estimation can use the exact
$m(v)$ of \eqref{eq:m-v}, the hint, and the hidden-error response correlations.

Breaking the 128-bit parameter set already contradicts the claimed security
of the submitted family at that level and supplies a practical UF-CMA forgery.

\section{Repair considerations}

Suppressing the constant response would remove the easiest sign sensor, but it
would not make \eqref{eq:accepted-mixture} secret-independent.  The remaining
response coordinates still share $r$ and retain secret-dependent correlations.
Likewise, using more precise scalar acceptance arithmetic or independently
resampling mask partitions cannot repair a joint-distribution error.

A repair must change the distributional argument.  Possible directions are:

\begin{enumerate}[leftmargin=1.7em]
  \item use rejection with the true vector mixture
  \[
   \tfrac12\prod_i f_i(z_i-v_i)+
   \tfrac12\prod_i f_i(z_i+v_i)
  \]
  in the denominator and prove a practical global rejection bound;
  \item adopt a standard bimodal rejection construction whose acceptance
        probability is derived for the complete vector, as in prior work on
        bimodal lattice signatures \cite{DDLL}; or
  \item remove the global bimodal sign and redesign the response and parameter
        set around a jointly sound unimodal rejection sampler.
\end{enumerate}

Any revised design should test joint moments and latent-variable distinguishers
in addition to checking scalar histograms.  Matching every one-dimensional
marginal is insufficient for zero knowledge or signing-key privacy.

\section{Conclusion}

BiT's rejection sampler proves and implements the wrong distributional object.
The scalar target is correct, but a sign shared by the full response survives
as a latent variable in the joint distribution.  The public response to the
known constant secret component exposes that variable with high confidence.
Posterior-weighted ring regression then recovers $s_0$, and public-key
compression both corrects the last uncertain coefficients and completes an
equivalent secret key.

The resulting attack is practical by a wide margin.  The main run required
200,000 signatures and 139 CPU-seconds on an older four-core machine while
using one thread.  Three of four tested keys forged at that query count, and
the fourth forged at 250,000.  Security analysis of bimodal rejection must
therefore preserve the global coupling throughout; coordinatewise marginal
arguments cannot justify a vector signature distribution.

\appendix
\section{Independent reproduction procedure}

The following procedure is sufficient to reproduce the attack without any
private experimental artifact.

\begin{enumerate}[leftmargin=1.7em]
  \item Obtain the official BiT candidate archive cited in \cite{BiT}, select
        the BiT-128 reference parameter set, and compile its key generation,
        signing, signature decoding, challenge expansion, and verification
        routines.  Confirm that public keys and signatures have lengths 1,048
        and 1,504 bytes.

  \item Generate one key pair.  Obtain 200,000 valid signatures on distinct
        messages.  Only the public key and signatures are inputs to the steps
        below.  Decode $(z_{1,0},z_{1,1},h,\widetilde c)$ and expand
        $\widetilde c$ to the weight-30 polynomial $c$ for each signature.

  \item For each of the 30 indices $i$ with $c_i=1$, read
        $x=z_{1,0,i}$.  Accumulate the log likelihood in \eqref{eq:llr}; if a
        numerator is zero set $w=-1$, and if a denominator is zero set $w=+1$.
        Otherwise set $w=\tanh(\Lambda/2)$.  As basic checks, decoded $x$ must
        lie in $[-7,7]$ and the running mean of $w^2$ should approach 0.889.

  \item Let
        $\zeta_k=\exp(\pi\mathrm{i}(2k+1)/256)$ and define
        $\widehat a_k=\sum_{j=0}^{255}a_j\zeta_k^j$.  For each of the three
        public $z_{1,1}$ polynomials, accumulate
        $B_{j,k}\mathrel{+}=w\overline{\widehat c_k}\widehat z_{j,k}$.
        Accumulate the common denominator
        $D_k\mathrel{+}=w^2|\widehat c_k|^2$.

  \item Compute $\widetilde s_j$ by the inverse negacyclic transform of
        $B_{j,k}/D_k$ and round every coefficient to the nearest member of
        $\{-1,0,1\}$.  No secret-key data is needed to assess this candidate:
        test whether every coefficient of $A_0\widetilde s$ admits a ternary
        error satisfying \eqref{eq:public-filter}.

  \item If the direct candidate fails, rank coefficients by
        $\bigl||\widetilde s_{j,i}|-1/2\bigr|$.  For the 24 lowest-confidence
        positions, pair the rounded value with its adjacent ternary value.
        Enumerate subsets in Gray-code order.  Precompute the contribution of
        each alternative to $A_0s_0$ and screen each subset on three public-key
        coefficients before applying the full check.

  \item For the passing $s_0$, choose $e'$ and calculate $b'_0$ as in
        \eqref{eq:completion}.  Form a signing key containing the submitted
        public-key components, $s_0,e',b'_0$, a fresh signing-randomness seed,
        and the hash of the public key.  Sign a message that was not among the
        oracle queries and submit the resulting 1,504-byte signature to the
        unchanged reference verifier.
\end{enumerate}

The transform accumulators are the only state that grows with the ring degree;
the query count affects running time but not attack memory.  Repeating the
experiment with a different key-generation seed provides an independent key
trial without changing any recovery parameter.

\begin{thebibliography}{9}

\bibitem{BiT}
X.~Lu, Y.~Yin, Y.~Liu, D.~Cai, D.~Jia, L.~Bi, Y.~Liu, B.~Shi, J.~He,
Y.~Shi, H.~Zhang, Y.~Liu, and Y.~Zhang.
\newblock \emph{BiT: Bimodal Triangular Distribution Based Lattice Signatures}.
\newblock Submission to the Next-generation Commercial Cryptographic
Algorithms Program, 2026.  Candidate archive:
\url{https://www.niccs.org.cn/niccs/Proposal/Public-Key%20Cryptographic%20Algorithms/Round%201%20candidates/BiT.zip}.

\bibitem{DDLL}
L.~Ducas, A.~Durmus, T.~Lepoint, and V.~Lyubashevsky.
\newblock Lattice signatures and bimodal Gaussians.
\newblock In \emph{Advances in Cryptology -- CRYPTO 2013}, volume 8042 of
LNCS, pages 40--56. Springer, 2013.
\newblock Full version: \url{https://eprint.iacr.org/2013/383}.

\bibitem{Lyu12}
V.~Lyubashevsky.
\newblock Lattice signatures without trapdoors.
\newblock In \emph{Advances in Cryptology -- EUROCRYPT 2012}, volume 7237 of
LNCS, pages 738--755. Springer, 2012.

\bibitem{Lyu09}
V.~Lyubashevsky.
\newblock Fiat--Shamir with aborts: Applications to lattice and
factoring-based signatures.
\newblock In \emph{Advances in Cryptology -- ASIACRYPT 2009}, volume 5912 of
LNCS, pages 598--616. Springer, 2009.

\bibitem{NGCC}
Institute of Commercial Cryptography Standards.
\newblock Submission requirements for public-key cryptographic algorithms.
\newblock Next-generation Commercial Cryptographic Algorithms Program, 2026.
\newblock \url{https://www.niccs.org.cn/niccs/Notice/1975896137741635584/lDop1mav.pdf}.

\end{thebibliography}

\end{document}
